% SPDX-License-Identifier: Apache-2.0
% covers: Vreteno
\datasheet{Vreteno}{A three-stage RV32IMAC core with machine, supervisor and user modes and an AXI host port.}

\dsfacts{\code{cpu/vreteno/src/core.rs}}{\code{vreteno32::core::Vreteno}}%
{\code{//docs:vreteno}, \code{//docs:noc}, \code{//docs:tapeout}}

\subsection*{Function}

\code{Vreteno} is a RISC-V core for the RV32I base set, the M
extension and the C extension's compressed instructions. It is one
process with three stages. The fetch stage reads the instruction at
the program counter into the instruction register, expanding a
compressed one into the thirty-two bit instruction it stands for. The
execute stage decodes that word, computes, and issues
loads and stores on the bus. The writeback stage extends a loaded
word, writes the register file and retires the instruction.

The core holds three memories: the register file, 32 words; the
instruction memory, \code{IMEM\_WORDS} or 1024 words; and a data RAM
of 64~KiB at \code{DRAM\_BASE}, \code{0x1\_0000}, the window a program
keeps its stack in (issue 1275). \code{Vreteno::with} loads a program
into the instruction memory. The data memory and every device are on
the bus, outside the core. Above
the instruction memory the fetch reads from the bus, into a buffer of
two words, and the words of the data memory and of DDR3 through an
instruction cache (issue 1021): sixteen kilobytes, 1024 lines of
four words, one way, looked up after translation and so indexed and
tagged by the physical address (issue 1290), in \code{ic\_data} and
\code{ic\_tag}. A hit fills both words of the buffer when the line
holds the next one, and the buffer shifts its second word into the
first at the clock edge, as the program counter moves into it
(issues 1187 and 1279).

The data RAM is on the core's own port and not on the bus: four lanes
of a byte, \code{dl0} to \code{dl3}, each written at one address and
read at another, as a block RAM is. A load in its window reads the
lanes into \code{dl\_word} at the edge and takes the word from there in
writeback, the cycle after it ran; a store writes the lanes under its
strobes in the cycle it runs. Whether an access is in the window comes
from the upper half of \code{rs1} and the carry out of the lower
halves' sum, without the full add.

The bus side is an AXI host written as hardware. The core issues a
burst on \code{issue}, sends a store's beat on \code{wbeat}, and hands
identifiers back on \code{release}. It receives \code{grant},
\code{done} and \code{rdata} from its tracker, an \code{AxiHost}. The
other inputs are \code{rst}, the external interrupt \code{irq}, the
timer's line \code{tirq} and the software interrupt \code{sirq}, which
the controller raises from \code{msip}; and the debug module's, the
halt and resume requests \code{haltreq} and \code{resumereq} and its
register access, \code{dbg\_regno}, \code{dbg\_wdata} and
\code{dbg\_we}, honoured in debug mode (issue 154). The outputs are
\code{halt}, \code{instr}, the word retiring, \code{wb}, a
\code{Writeback} of \code{done}, \code{rd} and \code{val}, \code{dbg},
debug mode, and \code{dbg\_rdata}, the register the module asked for.

Two additional inputs supply the machine timer count, \code{time}
(accessed by \code{rdtime}), and \code{seirq}, the supervisor external
interrupt line registered alongside MEIP as \code{mip.SEIP} (issue 1094).

Eleven ports join the core to its memory management unit, an
\code{Mmu8}, inside a \code{Hart} (issue 1014). The core tells the
unit \code{mmu\_satp}; \code{mmu\_prv}, the mode a data access is
checked as, which is \code{MPP} under \code{MPRV} (issue 1105);
\code{mmu\_sum} and \code{mmu\_mxr} from \code{mstatus};
\code{mmu\_flush}, raised for a write of
\code{satp} or an \code{sfence.vma}; and its two requests,
\code{ireq} for a fetch and \code{dreq} for a load or a store. The
unit answers on \code{ires} and \code{dres}, each a \code{Res} of a
physical address, a page fault or an access fault, a cycle after the
request. Its walker's reads come in on \code{ptw} as addresses, go out
on the core's own \code{issue}, and their words go back to the unit on
\code{pte}, so the hart needs no bus port of its own for the page
tables. The core believes an answer only at the second cycle of a
request, which is right both in simulation, where the unit sees a
request a step late, and in hardware.

\subsection*{Parameters}

\begin{center}\footnotesize
\begin{tabular}{@{}lp{0.7\textwidth}@{}}
\toprule
Parameter & Meaning \\
\midrule
\code{IW} & The width of the AXI identifier on \code{rdata}, \code{done},
\code{grant} and \code{release}, in bits. The tables use 2, as the
board and the runs do. \\
\code{DW} & The data RAM's words a lane, a power of two, 16384 by
default, which is 64~KiB. The netlist the documents simulate and the
layout maps uses 4, where the window's addresses wrap, since the cell
library has no RAM. The tables use the default. \\
\bottomrule
\end{tabular}
\end{center}

\subsection*{Register map}

The control registers, by CSR number. Software reaches them through
the six CSR instructions. They are not on the bus.

\begin{center}\footnotesize
\begin{tabular}{@{}llp{0.6\textwidth}@{}}
\toprule
Number & Register & What the core keeps \\
\midrule
\code{0x300} & \code{mstatus} & Bits 3 and 7 only; a write is masked
with \code{0x88}. \\
\code{0x304} & \code{mie} & Bits 11 and 7 only, the external and the
timer enable. \\
\code{0x305} & \code{mtvec} & The handler's address; a write clears
bits 1 and 0. \\
\code{0x340} & \code{mscratch} & Entire 32-bit register. \\
\code{0x341} & \code{mepc} & The trap's address; a write clears
bit 0. \\
\code{0x342} & \code{mcause} & Entire 32-bit register. \\
\code{0x343} & \code{mtval} & Entire 32-bit register. \\
\code{0x344} & \code{mip} & Bit 11 is stored: \code{irq} sets it and
a write keeps only that bit. Bit 7 reads as \code{tirq}. \\
\code{0x7c0} & \code{mhalt} & Nothing: it reads as zero. A write of an
odd value stops the core when that instruction retires. \\
\code{0x301} & \code{misa} & Nothing: it reads
\code{0x4000\_1104}, which is \code{MXL} of 1 and the letters
\code{C}, \code{I} and \code{M}. Writable and ignored, as the
specification allows. \\
\code{0xf11} & \code{mvendorid} & Nothing: zero, which means
unimplemented. Read only. \\
\code{0xf12} & \code{marchid} & Nothing: zero. Read only. \\
\code{0xf13} & \code{mimpid} & Nothing: zero. Read only. \\
\code{0xf14} & \code{mhartid} & Nothing: zero, this machine's one
hart. Read only. \\
\bottomrule
\end{tabular}
\end{center}

A CSR instruction that names any other number is an illegal word and
traps.

\dstables{Vreteno}

\subsection*{Behaviour}

\begin{itemize}
\item With \code{rst} high, the fetch counter goes to zero, the
instruction register is emptied, and the load wait and the sequencer
are cleared.
\item The fetch reads the instruction memory at bits 11 to 2 of the
program counter and at the word after. The halfword at the counter is
the upper half of its word when bit 1 is set. If its low two bits are
not both set it is a compressed instruction, expanded in the fetch;
otherwise the halfword after it is the instruction's upper half. The
counter moves on by two or four.
\item A fetch above the instruction memory goes to the bus when its
physical page is a device's, and to the cache when it is the data
memory's or DDR3's. The cache reads the line's tag and the word the
cycle after, at the address the fetch registered in \code{ic\_pa}. A
hit puts the word in the fetch buffer. A miss reads the line as one
burst of four beats, \code{len} 3, once nothing else of the core's is
on the bus, writes each beat into the line, makes the line valid on the
last, and looks the word up again. A refused beat ends the fetch with
the word's access fault, and leaves the line invalid.
\item \code{fence.i} and \code{rst} clear the cache's tags, one line a
cycle, and a lookup waits until they are clear. A store does not reach
the cache, so code written by stores runs after a \code{fence.i}.
\item \code{ir\_c} says the instruction in execute was compressed.
\code{jal} and \code{jalr} link the address two bytes on for a
compressed one and four for the rest.
\item An instruction in execute that reads the register the writeback
stage is about to write gets that value directly. When that value is a
load's, the instruction waits one cycle and reads the register file
instead.
\item A taken branch, \code{jal}, \code{jalr}, a trap and \code{mret}
redirect the fetch. The word fetched in that cycle is squashed, which
costs one bubble.
\item A load or store in the data RAM's window does not go to the bus.
A load's word comes from \code{dl\_word} the cycle after it ran, and
never sets \code{dev\_wait}; a store writes the lanes in the cycle it
runs and is not counted in \code{stores\_out}; \code{lr.w} and
\code{sc.w} act there as on the bus, and an AMO takes its word into
\code{wb\_dev} in its first cycle in writeback and stores to the lanes.
A load there clears \code{wb\_err}, since nothing refuses it.
\item Any other load or store whose address has any bit set above bit
11 goes on the bus as a burst of one beat, size 2, \code{INCR}. A store's beat
has the lanes it covers as its strobe.
\item A load or store waits for room on both \code{issue} and
\code{wbeat}, whatever its address.
\item A store is posted and the core runs on. A load sets
\code{dev\_wait} and sits in writeback, holding the pipeline, until
its answer lands in \code{wb\_dev}.
\item When a write response and a read beat arrive together, the
write's identifier goes back first and the read's waits a cycle. A
read's identifier goes back on its last beat, so once for a cache
fill's four. The core drops what it receives on \code{grant}.
\item \code{ecall} traps with cause 11 and \code{ebreak} with cause 3,
which takes its own address as \code{mtval}. A word the core does not
know traps with cause 2, and \code{mtval} takes the word. A halfword
that is no RV32C instruction expands to itself, which no thirty-two bit
instruction is, so it traps with cause 2 and \code{mtval} takes the
halfword. An unaligned load traps with cause 4 and an unaligned store
with cause 6, and \code{mtval} takes the address that was asked for.
Every other trap writes zero to \code{mtval}.
\item On a trap, \code{mepc} takes the instruction's address, bit 7 of
\code{mstatus} takes bit 3, bit 3 is cleared, and the fetch goes to
\code{mtvec}. \code{mret} sets bit 3 from bit 7, sets bit 7, and goes
to \code{mepc}.
\item The external interrupt is ready when bit 11 is set in both
\code{mie} and \code{mip}. The software interrupt is ready when bit 3
of \code{mie} is set and \code{sirq} is high, and the timer's when bit
7 of \code{mie} is set and \code{tirq} is high. Any of the three is
taken when bit 3 of \code{mstatus} is set, in place of the instruction
in execute. The external interrupt, cause \code{0x8000\_000b}, wins
over the software interrupt, cause \code{0x8000\_0003}, which wins over
the timer's, cause \code{0x8000\_0007}.
\item A multiply or divide starts the sequencer and stalls in execute.
A multiply is one step and a division thirty-two. \code{//docs:vreteno}
states that a multiply takes three cycles and a division thirty-three.
An interrupt cancels the sequencer, and the instruction starts it
again after the handler returns.
\item The sequencer does not start while a device load in writeback is
still waiting for its answer.
\item A write of an odd value to \code{mhalt} stops the fetch and
parks the program counter one word past that instruction, which
retires first. \code{halt} rises as it retires, and stays high.
\item \code{stall}, \code{int\_take} and \code{redirect} are wires kept
as fields, so the trace and the netlist show them.
\end{itemize}

\subsection*{Verification}

\begin{itemize}
\item \code{//cpu/vreteno:lockstep\_test} runs the core with the router, the
data memory, the timer and the serial port against the reference model
in \code{model.rs}. After every cycle the program counter, registers
1 to 31, the eight control registers that hold something and the halt
must agree. At the
halt the data memory, the timer's compare and the serial port's bytes
must agree. Its tests are \code{demo\_program},
\code{a\_multiply\_right\_after\_a\_load},
\code{compressed\_instructions}, \code{an\_unaligned\_access\_traps},
\code{the\_data\_ram\_on\_the\_cores\_own\_port}, a store and a load of
one word at once, a load through a pointer just loaded, bytes and
halves, both arms of the window's compare, \code{lr}/\code{sc} and an
AMO there, and \code{random\_programs}, which
runs 64 seeds of 200 instructions each with both lengths mixed and
counts that the mix is there.
\item \code{//cpu/vreteno:isa\_test} compares the core's expander with
the decoder's on every compressed halfword, and
\code{//cpu/vreteno:compressed\_test} checks the decoder against the
GNU disassembler on all 49\,152 of them.
\item \code{//cpu/vreteno:hello\_test} (\code{rust\_runs\_on\_vreteno})
and \code{//cpu/vreteno:hello\_cc\_test}
(\code{cpp\_runs\_on\_vreteno}) run compiled programs on the machine
of \code{run.rs} and check what the serial line says and that the core
halted.
\item \code{//cpu/vreteno:board\_test} and \code{//soc:demo\_test} run
the core inside \code{Board} and on the lattice.
\item \code{//cpu/vreteno:icache\_test} runs issue 1009's loop of three
instructions from the instruction memory and from the data memory,
through the cache, at 1.334 and 3.349 cycles an instruction counted
from the reset. The second includes the 1024 cycles in which the cache
clears its tags, which the loop waits for, since the cache is sixteen
kilobytes (issue 1290); it was 3.093 at four kilobytes, 4.093 before
one cache hit filled both words of the buffer and the shift moved to
the edge (issues 1187 and 1279), and 14.655 from the data memory
before the cache. It also runs the same loop across two lines, and a
routine rewritten by a store and called again after
\code{fence.i}, which runs the word stored.
\item The build simulates the netlist of \code{Hart<2>}, the core and its
unit, under the entity name \code{vreteno}, with the demonstration
program in the core's instruction memory, against the trace of
\code{//cpu/vreteno:vreteno} under nvc and Verilator:
\code{//docs:vreteno\_sim\_vreteno}\allowbreak\code{\_vreteno\_tb\_test} and
\code{//docs:vreteno\_vsim\_vreteno\_test}.
\item \code{//cpu/vreteno:vreteno\_synth} and
\code{//cpu/vreteno:vreteno\_pnr} synthesise and place the core on an
\code{xc7a200tfbg484-2}; both are manual. \code{//docs:vreteno}
reports 2151 LUTs as logic, 48 LUTs as distributed RAM, 489
flip-flops, two block RAM tiles and four DSP blocks after synthesis. It
reports that the routed core meets a 10~ns clock with 0.296~ns of
setup slack.
\item \code{//cpu/vreteno/asic:vreteno\_cells} maps the core onto the
Nangate45 cells and builds with everything.
\code{//cpu/vreteno/asic:}\allowbreak\code{vreteno\_pnr} places and routes it, and is
manual.
\end{itemize}

\subsection*{Used by}

\code{Board}, as \code{cpu}. The demonstration
\code{//cpu/vreteno:vreteno}, the lockstep test, and \code{run.rs},
which the \code{hello} and \code{hello\_cc} runs call. The system in
\code{soc}, where the core is the host at corner 0,0 of the lattice.

\subsection*{Limits}

The boot memory is 1024 words, 4096 bytes, and nothing loads it at
run time: a netlist gets its contents through \code{init}. A program
counter at or above \code{IMEM\_BYTES}, the boot memory's end, is
fetched from the bus instead (issue 134), so a program larger than the
boot memory runs from wherever it was loaded, the DDR3 memory on the
board.

Every load and store goes to the bus, the boot memory included, which
answers reads and refuses writes (issue 268), except those in the data
RAM's window. Nothing but the core reaches the data RAM: not the debug
module's system bus access, Razboj or a DMA engine, so a debugger
reads it through the core. A fetch from the window is refused.

A trap handler must start at a whole word: \code{mtvec} keeps no
bits 1 and 0. The core has machine mode only and the fourteen control
registers above, in three kinds: eight it keeps, \code{mhalt}, which
is its own, and five that say what the machine is and hold nothing. A
write to one of the four whose number begins \code{0xf} is an illegal
instruction, which is what read only means in the specification;
\code{misa} is writable and ignores what is written. The two counters answer as well:
\code{mcycle} at \code{0xb00} with its high half at \code{0xb80},
and \code{minstret} at \code{0xb02} with \code{0xb82}, each 64 bits
and each writable. \code{mcycle} counts while the core runs and stops
when the halt stops it; \code{minstret} counts what retires. They are
the one part of the core the lockstep test leaves alone, for the
reason \code{//docs:vreteno} gives, and
\code{//cpu/vreteno:counters\_test} checks them instead. It has one load in
flight at a time. On the board, \code{irq} is the interrupt
controller's line.
